
% TODO: in either this chapter or the next, mention keys() as "a relevant thing
% to do" when you're trying to get your head around what a big data set
% contains. (If it's a list (which you can tell with type(s)), you can x[0] it.
% If it's a dictionary, you can keys(x) it and then x['whatever'] it.)

\chapter{JSON (1 of 2)}
\label{ch:JSON1}

JSON (pronounced like the name ``Jason'') is a simple, human-readable text file
format that is very commonly used for information storage and interchange. (It
stands for ``JavaScript Object Notation,'' but don't be fooled: it doesn't have
much to do with the JavaScript language except historically. \textit{Any}
language, including Python, can read/write JSON data.)

In this chapter, we'll learn about how to read and navigate through JSON data.
First, though, a word about a couple of Python essentials I omitted from DATA
101 for brevity and simplicity.

\section{Lists and dictionaries}

\index{series@\texttt{Series} (Pandas)}
\index{ndarray@\texttt{ndarray} (NumPy)}
\index{associative array}
\index{key-value pair}
\index{index}

Last semester, we spent a lot of time learning how to use NumPy
\texttt{ndarray}s and Pandas \texttt{Series}es. These are the preferred
implementations of arrays, and associative arrays, respectively, in the Python
Data Science ecosystem. As you know, a NumPy array holds a numbered sequence of
items (indexed starting from 0), and a \texttt{Series} holds a set of key-value
pairs.

For any serious data work, you want to use those tools, because they are
lightning-fast, super-efficient, feature-rich, and optimized to the task.

\index{plain-ol' Python}

As it turns out, plain-ol' Python has stripped-down, feature-poor, but
easy-to-type versions of these data structures as well, which are baked in to
JSON and also important to know about in their own right. They're called
\textbf{lists} and \textbf{dictionaries}.

\subsection{Lists}

If, instead of writing this:

\begin{Verbatim}[fontsize=\small,samepage=true,frame=single,framesep=3mm]
crew = np.array(["Ed", "Kelly", "Alara", "Bortus", "John", "Claire"])
\end{Verbatim}

you simply wrote this:

\begin{Verbatim}[fontsize=\small,samepage=true,frame=single,framesep=3mm]
crew = ["Ed", "Kelly", "Alara", "Bortus", "John", "Claire"]
\end{Verbatim}

\index{list}

then you get a plain-ol \textbf{list} instead of a NumPy array. It isn't nearly
as fast, nor can it store nearly as many items, and you can't do broadcasting
operations and such with it, but it does the basic job. For example:

\index{len@\texttt{len()} function!for lists}

\begin{Verbatim}[fontsize=\small,samepage=true,frame=single,framesep=3mm]
print("Captain {} initiating roll call for {} officers:".format(
    crew[0], len(crew)))
for member in crew[1:len(crew)]:
    print("{} reporting for duty, sir!".format(member))
print(crew[0])
\end{Verbatim}

\vspace{-.2in}

\begin{Verbatim}[fontsize=\small,samepage=true,frame=leftline,framesep=5mm,framerule=1mm]
Captain Ed initiating roll call for 6 officers:
Kelly reporting for duty, sir!
Alara reporting for duty, sir!
Bortus reporting for duty, sir!
John reporting for duty, sir!
Claire reporting for duty, sir!
\end{Verbatim}

Working with small lists of items is very, very similar to working with NumPy
arrays, so I'm sure you'll pick this right up.

Incidentally, we've actually used lists several times already, without knowing
it. For instance, this very line of code:

\begin{Verbatim}[fontsize=\small,samepage=true,frame=single,framesep=3mm]
good = np.array(["Janet", "Chidi", "Eleanor", "Michael", "Tahani"])
\end{Verbatim}

is actually first creating a \textit{list}, and then passing that list as an
argument to \texttt{np.array()}. And this code:

\begin{Verbatim}[fontsize=\small,samepage=true,frame=single,framesep=3mm]
alter_egos = pd.Series(['Hulk','Spidey','Iron Man','Thor'],
    index=['Bruce','Peter','Tony','Thor'])
\end{Verbatim}

creates two separate lists, one for the data and one for the index, and passes
those to \texttt{pd.Series()} to build a \texttt{Series} out of them. Lists are
super-lightweight and super-common.

\subsection{Dictionaries}

\index{series@\texttt{Series} (Pandas)}
\index{dictionary}
\index{key-value pair}

As mentioned in section~\ref{mentionDictionary}
(p.~\pageref{mentionDictionary}), \textbf{dictionaries} are like stripped-down
Pandas \texttt{Series}es, in the same way that lists are like stripped-down
NumPy arrays. Dictionaries contain key-value pairs, accessed in the same sort
of way (with the boxie \texttt{[]} syntax). You create a dictionary by
using the syntax from p.~\pageref{mentionDictionary2}:

\begin{Verbatim}[fontsize=\small,samepage=true,frame=single,framesep=3mm]
alter_egos = { 'Bruce':'Batman', 'Dick':'Robin', 'Diana':'Wonder Woman',
    'Clark':'Superman' }
\end{Verbatim}

Then, just like with \texttt{Series}es, you can retrieve the value for a key,
change the value for a key, add a new key-value pair, and so forth:

\index{len@\texttt{len()} function!for dictionaries}

\begin{Verbatim}[fontsize=\footnotesize,samepage=true,frame=single,framesep=3mm]
caped_crusaders = [ alter_egos['Bruce'], alter_egos['Dick'] ]
alter_egos['Clark'] = 'Man of Steel'
alter_egos['Victor'] = 'Cyborg'

print("There are now {} heroes, {} of which are the caped crusaders.".format(
    len(alter_egos), len(caped_crusaders)))
print("The dynamic duo is: {} and {}!".format(caped_crusaders[0],
    caped_crusaders[1]))
for key, value in alter_egos.items():
    print("Psssst, {} is really {}...".format(key,value))
\end{Verbatim}
\vspace{-.2in}

\begin{Verbatim}[fontsize=\small,samepage=true,frame=leftline,framesep=5mm,framerule=1mm]
There are now 5 heroes, 2 of which are the caped crusaders.
The dynamic duo is: Batman and Robin!
Psssst, Bruce is really Batman...
Psssst, Dick is really Robin...
Psssst, Diana is really Wonder Woman...
Psssst, Clark is really Man of Steel...
Psssst, Victor is really Cyborg...
\end{Verbatim}

\subsection{Nesting...and nesting, and nesting...}

\index{nested!data structures}

All that was easy so far. The only thing left to tell you -- and man, it can
pack a whammy when you first see it -- is that both lists and dictionaries can
be \textbf{nested} indefinitely. This means that one of a list's elements might
well be another list, or a dictionary. And the value for one of a dictionary's
keys might well be a list, or another dictionary.

Sometimes you have to stare hard at the printout of a nested structure to
figure out how to unpack it. Here are some examples:

\begin{Verbatim}[fontsize=\small,samepage=true,frame=single,framesep=3mm]
a = [ 7, 1, [9, 5], 6 ]
\end{Verbatim}

\index{[]@\texttt{[]} (boxies)}
\index{boxies (square brackets)}

The \texttt{a} variable is a list (you can tell because it starts with a boxie)
and one of its elements happens to be another list (ditto). Stare hard at that,
and then stare hard at this output, and make sure you see why it's so:

\begin{Verbatim}[fontsize=\small,samepage=true,frame=single,framesep=3mm]
print("a[1] is {}".format(a[1]))
print("len(a) is {}".format(len(a)))
print("a[3] is {}".format(a[3]))
print("len(a[2]) is {}".format(len(a[2])))
print("a[2][1] is {}".format(a[2][1]))
\end{Verbatim}
\vspace{-.2in}

\begin{Verbatim}[fontsize=\small,samepage=true,frame=single,framesep=3mm]
a[1] is 1
len(a) is 4
a[3] is 6
len(a[2]) is 2
a[2][1] is 5
\end{Verbatim}

The first one was a warmup. But the second output may have surprised you:
aren't there \textit{five} numbers in that list, not four? The answer is
\textit{no}, there are exactly four:

\begin{compactitem}
\item Item \#0 is the number 7.
\item Item \#1 is the number 1.
\item Item \#2 is a list, which has two elements: 9 and 5.
\item Item \#3 is the number 6.
\end{compactitem}

Once you grasp that, you can probably understand the other outputs from that
example as well. Of particular note is the last one: \texttt{a[2][1]} means
``get element \#2 from \texttt{a}, and since it is itself a list, get element
\#1 from \textit{that}.'' And since all lists are numbered starting with 0,
that's effectively the second element of the third element of \texttt{a}, which
is the number 5.

Next one:

\index{apple}
\index{banana (fruit)}
\index{cherry}
\index{bing cherry}
\index{rainier cherry}

\begin{Verbatim}[fontsize=\small,samepage=true,frame=leftline,framesep=5mm,framerule=1mm]
b = { "apple": 54, "banana": [14, 67], "cherry": { "bing": [ 13, 17,
21], "rainier": { "sweet": "yes", "juicy": "yes" } } }
\end{Verbatim}

\index{\{\}@\texttt{\{\}} (curlies)}
\index{curlies (curly braces)}

The \texttt{b} variable is a dictionary (you can tell because it starts with a
curly) and it contains several key-value pairs. How many, you ask?

\begin{Verbatim}[fontsize=\small,samepage=true,frame=single,framesep=3mm]
print(len(b))
\end{Verbatim}
\vspace{-.2in}

\begin{Verbatim}[fontsize=\small,samepage=true,frame=leftline,framesep=5mm,framerule=1mm]
3
\end{Verbatim}

The three key-value pairs are:

\begin{compactitem}
\item \texttt{"apple"}, whose value is the number 54.
\item \texttt{"banana"}, whose value is a list, whose elements are 14 and 67.
\item \texttt{"cherry"}, whose value is another dictionary. \textit{That}
dictionary has two key-value pairs:
    \begin{compactitem}
    \item \texttt{"bing"}, whose value is a list with 13, 17, and 21.
    \item \texttt{"rainier"}, whose value is yet another dictionary, with keys
\texttt{"sweet"} and \texttt{"juicy"}.
    \end{compactitem}
\end{compactitem}

Now test yourself and make sure you can predict the right output from these
lines of code:

\begin{Verbatim}[fontsize=\small,samepage=true,frame=single,framesep=3mm]
print("b['apple'] is {}".format(b['apple']))
print("b['banana'] is {}".format(b['banana']))
print("len(b['banana']) is {}".format(len(b['banana'])))
print("len(b['cherry']) is {}".format(len(b['cherry'])))
print("len(b['cherry']['bing']) is {}".format(len(b['cherry']['bing'])))
print("b['cherry']['bing'][0]) is {}".format(b['cherry']['bing'][0]))
print("b['cherry']['rainier']['sweet']) is {}".format(
    b['cherry']['rainier']['sweet']))
\end{Verbatim}
\vspace{-.2in}

\begin{Verbatim}[fontsize=\small,samepage=true,frame=leftline,framesep=5mm,framerule=1mm]
b['apple'] is 54
b['banana'] is [14, 67]
len(b['banana']) is 2
len(b['cherry']) is 2
len(b['cherry']['bing']) is 3
b['cherry']['bing'][0]) is 13
b['cherry']['rainier']['sweet']) is yes
\end{Verbatim}

\index{len@\texttt{len()} function!for lists}

Remember, \texttt{len()} plays several roles. When called on a list or array,
it gives the number of elements; on a dictionary or \texttt{Series}, the number
of key-value pairs; on a string, the number of characters; and on a
\texttt{DataFrame}, the number of rows. That's why \texttt{len(b['banana'])}
and \texttt{len(b['cherry'])} both return 2, even though they're very different
types of things.

\index{type@\texttt{type()}}

Note this general rule, too: every time the thing you have is a \textit{list},
you must follow it with an \textit{integer} in boxies. Every time the thing you
have is a \textit{dictionary}, you must follow it with a \textit{key} (often a
string) in boxies. Always.

And now's a good time to mention how to find out what type of thing something
\textit{is} -- the \texttt{type()} function:

\begin{Verbatim}[fontsize=\small,samepage=true,frame=single,framesep=3mm]
print("b['apple'] is a {}".format(type(b['apple'])))
print("b['banana'] is a {}".format(type(b['banana'])))
print("b['cherry'] is a {}".format(type(b['cherry'])))
print("b['cherry']['rainier']['juicy'] is a {}".format(
    type(b['cherry']['rainier']['juicy'])))
\end{Verbatim}
\vspace{-.2in}

\begin{Verbatim}[fontsize=\small,samepage=true,frame=leftline,framesep=5mm,framerule=1mm]
b['apple'] is a int
b['banana'] is a list
b['cherry'] is a dict
b['cherry']['rainier']['juicy'] is a str
\end{Verbatim}

This is super useful when plunging deep through nested data sets.

One last example:

\begin{Verbatim}[fontsize=\small,samepage=true,frame=single,framesep=3mm]
c = [ { 'x': 5 }, [ 14, 'bob', [ 11, { 'duck': 'donuts', 'five':
'guys' }, [ 17, 6 ] ] ] ]
\end{Verbatim}

\index{[]@\texttt{[]} (boxies)}
\index{boxies (square brackets)}

Can you visually line up all those boxies? If so, then you can probably get
these answers correct:

\begin{Verbatim}[fontsize=\small,samepage=true,frame=single,framesep=3mm]
print("len(c) is {}".format(len(c)))
print("c[0]['x'] is {}".format(c[0]['x']))
print("c[1][1] is {}".format(c[1][1]))
print("c[1][2][1]['five'] is {}".format(c[1][2][1]['five']))
print("c[1][2][2][1] is {}".format(c[1][2][2][1]))
print(\end{Verbatim}
\vspace{-.2in}

\begin{Verbatim}[fontsize=\small,samepage=true,frame=leftline,framesep=5mm,framerule=1mm]
len(c) is 2
c[0]['x'] is 5
c[1][1] is bob
c[1][2][1]['five'] is guys
c[1][2][2][1] is 6
\end{Verbatim}



\section{JSON}

\index{JSON}
\index{XML}

Okay, now about JSON proper.

The really nice things about JSON are (1) it's a simple format, and relatively
easy on the eyes (unlike binary formats, or XML), and (2) it structures
information in nested list-ish, dictionary-ish, string-ish type
elements that are exactly like what we covered in the previous section. A JSON
``document'' is, in fact, \textit{essentially a Python list or dictionary with
embedded contents.}

Here's an example file in JSON format:

\begin{Verbatim}[fontsize=\footnotesize,samepage=true,frame=lines,framesep=3mm]
{"name": "Luke Skywalker", "rank": "Jedi", "lightsaber color": "blue",
"episodes": ["IV", "V", "VI", "VIII", "IX"], "relations": {"Leia": "sister",
"Han": "friend", "Artoo": "servant", "Threepio": "servant", "Rey": "student"}}
\end{Verbatim}

Its structure is instantly apparent: a dictionary whose values are strings,
lists, and other dictionaries, as appropriate.

Probably the best thing about JSON is that once we read and parse a JSON file
or other information source, we have a bona fide Python variable that we can
treat in the ordinary way.

In Python, we have the \texttt{json} package, which you can import like any
other:

\begin{Verbatim}[fontsize=\small,samepage=true,frame=single,framesep=3mm]
import json
\end{Verbatim}

There are four functions that matter:
\begin{compactitem}

\item \texttt{json.dump}(\textit{someVariable}, \textit{aFile}) - Dump the
Python variable to a file in JSON text format.

\item \texttt{json.dumps}(\textit{someVariable}) - Dump the Python
variable to a string in JSON text format (and return it).

\item \texttt{json.load}(\textit{aFile}) - Build and return a Python
variable from the contents of the JSON text file.

\label{json.loads}
\item \texttt{json.loads}(\textit{someString}) - Build and return a Python
variable from the JSON text string.

\end{compactitem}

\index{json file@\texttt{.json} file}

So, for instance, to load a file \texttt{info.json} into a Python variable
called ``\texttt{my\_stuff}'':

\begin{Verbatim}[fontsize=\small,samepage=true,frame=single,framesep=3mm]
with open('info.json','r') as f:
    my_stuff = json.load(f)
\end{Verbatim}

We could parse and access the Luke Skywalker file just like this:

\begin{Verbatim}[fontsize=\small,samepage=true,frame=single,framesep=3mm]
with open('luke.json','r') as f:
    luke = json.load(f)
print(luke['lightsaber color'])
blue

print(luke['episodes'][2])
VI
\end{Verbatim}

\normalsize
Simple as that!

\subsection{``Pretty printing''}

\index{pretty printing@``pretty printing''}
\index{wall of text}

One thing I think you'll find useful is the ability to ``\textbf{pretty
print}'' a nested Python structure, like one you read from a JSON file. Instead
of a monolithic wall of text, like many \texttt{.json} files look out of the
box, a pretty printer attempts to format, indent, and put linebreaks in, so
that you can more readily perceive its structure and figure out how to pick out
the parts you want.

A good package for this is the \texttt{pprint} package. First, you have to do
the import:

\index{import@\texttt{import}}

\begin{Verbatim}[fontsize=\small,samepage=true,frame=single,framesep=3mm]
from pprint import pprint
\end{Verbatim}

\index{print()@\texttt{print()} function}
\index{pprint()@\texttt{pprint()} function}

From then on, instead of \texttt{print()}ing a variable, you can
\texttt{pprint()} it. I think you'll agree, the latter is easier to make sense
of:

\begin{Verbatim}[fontsize=\small,samepage=true,frame=single,framesep=3mm]
print(luke)
\end{Verbatim}
\vspace{-.2in}

\begin{Verbatim}[fontsize=\footnotesize,samepage=true,frame=leftline,framesep=5mm,framerule=1mm]
{"name": "Luke Skywalker", "rank": "Jedi", "lightsaber color": "blue",
"episodes": ["IV", "V", "VI", "VIII", "IX"], "relations": {"Leia": "sister",
"Han": "friend", "Artoo": "servant", "Threepio": "servant", "Rey": "student"}}
\end{Verbatim}

\begin{Verbatim}[fontsize=\small,samepage=true,frame=single,framesep=3mm]
pprint(luke)
\end{Verbatim}
\vspace{-.2in}

\begin{Verbatim}[fontsize=\small,samepage=true,frame=leftline,framesep=5mm,framerule=1mm]
{'episodes': ['IV', 'V', 'VI', 'VIII', 'IX'],
 'lightsaber color': 'blue',
 'name': 'Luke Skywalker',
 'rank': 'Jedi',
 'relations': {'Artoo': 'servant',
               'Han': 'friend',
               'Leia': 'sister',
               'Rey': 'student',
               'Threepio': 'student'}}
\end{Verbatim}

As a bonus, it even puts dictionary keys in alphabetical order, so you can find
them more easily.


\section{``Hierarchical'' vs. ``flat'' data}

\index{flat data}
\index{hierarchical data}

So you see that the JSON file format is essentially a textual representation of
an Python variable with arbitrarily nested inner structure. This variable can
be a list or a dictionary, and have embedded lists and dictionaries to any
level of nesting. Here's another example:

\index{courses (JSON file)}
\label{coursesJson}

\begin{Verbatim}[fontsize=\small,samepage=true,frame=lines,framesep=3mm]
[{"prefix": "DATA", "num": 219, "creds": 3, "instructor": "Davies",
"meetings": [{"time": {"day": "TR", "time": "12:30pm"}, "location": {"bldg":
"Trinkle", "room": "B6"}}], "roster": [{"first": "Lady", "last": "Gaga",
"email": "lgaga@umw.edu"}, {"first": "The", "last": "Rock", "email":
"trock@umw.edu"}, ...
\end{Verbatim}

When you disentangle this, you'll realize that what you have is:

\vspace{.1in}

A list, each element of which is:
\vspace{-.1in}
\begin{compactitem}
\item a dictionary, with key/value pairs:
        \begin{compactitem}
        \item \texttt{prefix}: a string
        \item \texttt{num}: an int
        \item \texttt{creds}: an int
        \item \texttt{instructor}: a string
        \item \texttt{meetings}: a list, each element of which is:
            \begin{compactitem}
            \item a dictionary, with key/value pairs:
                \begin{compactitem}
                \item \texttt{time}: a dictionary, with key/value pairs:

\quad\quad \texttt{day}: a string

\quad\quad \texttt{time}: a string

                \item \texttt{location}: a dictionary, with key/value pairs:

\quad\quad \texttt{bldg}: a string

\quad\quad \texttt{room}: a string

                \end{compactitem}
            \end{compactitem}
        \item \texttt{roster}: a list, each element of which is:
            \begin{compactitem}
            \item a dictionary, with key/value pairs:
                \begin{compactitem}
                \item \texttt{first}: a string
                \item \texttt{last}: a string
                \item \texttt{email}: a string
                \end{compactitem}
            \end{compactitem}
        \item ...and possibly other stuff...
        \end{compactitem}
\end{compactitem}

\normalsize

\pagebreak

This kind of structure is called \textbf{hierarchical} data. This is in
contrast to so-called \textbf{flat} data, like a CSV:
\index{CSV file}

\begin{Verbatim}[fontsize=\normalsize,samepage=true,frame=lines,framesep=3mm]
Prefix,Num,Instructor,Creds,Bldg,Room,Day,Time
DATA,219,Davies,3,Trinkle,B6,TR,12:30pm
ENGL,286,Rigsby,3,Combs,112,MWF,11am
CPSC,305,Finlayson,4,Trinkle,B52,MWF,8am
\end{Verbatim}

It should be clear that hierarchical data is intrinsically \textit{more
expressive} than flat data. Rather than each data record being a strict,
identical-length sequence of values, the whole nested structure communicates
order and relationships between subcomponents, like an outline.

It's not even always possible, let alone straightforward, to convert from a
hierarchical format to a flat format. Consider the JSON file on
p.~\pageref{coursesJson}: it permits any number of \texttt{meetings} for a
single course, each of which has embedded structure. In this case, the
structure for each meeting is regular, since it's apparently composed of a
fixed number of fields with unique names. (This won't always be the case, and
if it's not, we would \textit{lose information} in collapsing it to a flat
format.) Even so, though, it's not clear what to do with the multiple meeting
times/locations for a course. Do we make each \texttt{Bldg/Room/Day/Time} group
into its own set of headings in the flat format? This brings with it its own
problems, such as:

\begin{compactenum}
\item How do we associate each group of buildings/rooms/days/times together,
since they're now all individual top-level fields? (Naming conventions seem
like the only solution, and it's awkward.)
\item How do we know the maximum number of such meetings, so as to specify the
width of the table?
\item What do we do for rows which don't have the maximum number of meetings?
Leave the others blank?
\item How do we perform standard queries like ``find all courses that have at
least one meeting in Trinkle?'' (We have to split this up into an awkward query
that separately queries each of the columns and unions the answer.)
\end{compactenum}

There are better ways to organize such data, such as the \textbf{relational
model} that has been ubiquitous in database schema design since Ted Codd
invented it in the 1960's. It allows us to ``normalize'' such data sets and
represent them in multiple, interrelated tables. We'll take a look at this
when we look at SQLite later in the course. Some irregular structures defy
even the relational model, however, which is why some newer database paradigms
(like NoSQL databases) have been developed to get around its limitations.

One important fact to grasp, however, is that at the present state of
technology, \textit{almost all machine learning algorithms are based on a
\textbf{flat} data structure}. There are some quite advanced, cutting-edge
exceptions to this, but really the field of mining hierarchical (and other
non-flat) data formats is still in its infancy. Lots of opportunity to make a
name for yourself here!
